% ─────────────────────────────────────────────────────────────────────────────
%  Guía Módulo 02 — Cadenas de Markov: Ejercicio Genética Poblacional
%  Autor: P. Hurtado · 2026
% ─────────────────────────────────────────────────────────────────────────────
\documentclass[10pt,a4paper]{article}

\usepackage[T1]{fontenc}
\usepackage[utf8]{inputenc}
\usepackage{lmodern}
\usepackage[a4paper, top=1.8cm, bottom=2.0cm, left=1.5cm, right=1.5cm]{geometry}
\usepackage[spanish,es-noshorthands]{babel}
\usepackage{amsmath,amssymb,mathtools,amsthm}
\usepackage{xcolor}
\usepackage[most]{tcolorbox}
\usepackage{tikz}
\usepackage{pgfplots}
\pgfplotsset{compat=1.18}
\usepackage{microtype}
\usepackage{needspace}
\usepackage{parskip}
\usepackage{fancyhdr}
\usepackage{hyperref}
\usepackage{booktabs}
\usepackage{array}
\usepackage{caption}
\usepackage{enumitem}

\usetikzlibrary{arrows.meta, positioning, calc, shapes.geometric, fit, backgrounds, decorations.pathreplacing}

\setlength{\parskip}{4pt}
\setlength{\parindent}{0pt}
\setlength{\headheight}{19pt}
\allowdisplaybreaks

% ─── Colores ─────────────────────────────────────────────────────────────────
\definecolor{sectionbg}{rgb}{0.16,0.26,0.52}
\definecolor{subsecbg}{rgb}{0.38,0.50,0.75}
\definecolor{codeblue}{rgb}{0.0,0.28,0.68}
\definecolor{notebg}{rgb}{1.0,0.97,0.87}
\definecolor{noteborder}{rgb}{0.80,0.65,0.25}
\definecolor{exbg}{rgb}{0.96,0.96,0.96}
\definecolor{exborder}{rgb}{0.55,0.55,0.55}

\hypersetup{colorlinks=true, linkcolor=codeblue, citecolor=codeblue, urlcolor=codeblue}

% ─── Cabeceras ───────────────────────────────────────────────────────────────
\pagestyle{fancy}
\fancyhf{}
\fancyhead[L]{\footnotesize\sffamily\color{sectionbg}Gu\'ia M\'odulo 02 --- Modelos Estoc\'asticos}
\fancyhead[R]{\footnotesize\sffamily\color{sectionbg}2026}
\fancyfoot[C]{\footnotesize\sffamily\thepage}
\renewcommand{\headrulewidth}{0.4pt}
\renewcommand{\headrule}{\hbox to\headwidth{\color{subsecbg}\leaders\hrule height \headrulewidth\hfill}}

% ─── tcolorbox styles ────────────────────────────────────────────────────────
\tcbset{
  notebox/.style={
    breakable,
    colback=notebg, colframe=noteborder, boxrule=0.4pt, arc=3pt,
    left=5pt, right=5pt, top=4pt, bottom=4pt, fontupper=\footnotesize\itshape,
  },
}

% ─── Macros ──────────────────────────────────────────────────────────────────
\newcommand{\Prob}[1]{\mathbb{P}\!\left(#1\right)}
\newcommand{\E}[1]{\mathbb{E}\!\left[#1\right]}
\newcommand{\Pmat}{\mathbf{P}}
\newcommand{\pij}[2]{p_{\,#1#2}}

% ─── Encabezado de parte ─────────────────────────────────────────────────────
\newcommand{\parteheader}[1]{%
  \vspace{8pt}\needspace{4cm}\noindent
  \begin{tcolorbox}[
    colback=subsecbg!18, colframe=subsecbg!55,
    sharp corners, boxrule=0.4pt,
    left=6pt, right=6pt, top=3pt, bottom=3pt,
    fontupper=\small\bfseries\sffamily,
  ]
  #1
  \end{tcolorbox}\vspace{2pt}%
}

% ─── Entorno enunciado ───────────────────────────────────────────────────────
\newenvironment{enunciado}[1]{%
  \begin{tcolorbox}[
    breakable,
    colback=sectionbg!8, colframe=sectionbg,
    boxrule=0.8pt, arc=3pt, left=8pt, right=8pt, top=6pt, bottom=6pt,
    fonttitle=\normalsize\bfseries\sffamily, coltitle=white, colbacktitle=sectionbg, titlerule=0pt,
    title={#1},
  ]%
}{\end{tcolorbox}}

\newenvironment{nota}{\begin{tcolorbox}[notebox]}{\end{tcolorbox}}

% ─────────────────────────────────────────────────────────────────────────────
\begin{document}
% ─────────────────────────────────────────────────────────────────────────────

\begin{tcolorbox}[
    colback=sectionbg, colupper=white, colframe=sectionbg,
    sharp corners, left=12pt, right=12pt, top=10pt, bottom=10pt,
  ]
  {\Large\bfseries\sffamily Cadenas de Markov: Gen\'etica Poblacional}\\[4pt]
  {\normalsize\sffamily Gu\'ia M\'odulo 02 --- Modelos Estoc\'asticos}\\[6pt]
  {\small\sffamily\color{white!80!sectionbg}
    Espacio de estados \textperiodcentered\ Probabilidades de transici\'on
    \textperiodcentered\ Clasificaci\'on \textperiodcentered\ Absorci\'on}\\[6pt]
  {\footnotesize\sffamily\color{white!85!sectionbg}\itshape
    Por P.~Hurtado \textperiodcentered\ 2026}
\end{tcolorbox}

\vspace{6pt}

% ═══════════════════════════════════════════════════════════════════════
\begin{enunciado}{Problema 1 \textperiodcentered\ Modelo de Wright--Fisher (2 individuos)}
  Considere una poblaci\'on con s\'olo dos individuos que en cada generaci\'on producen
  dos nuevos individuos y despu\'es mueren (el tama\~no de la poblaci\'on es, por tanto,
  constante igual a dos). Cada individuo porta dos genes que determinan el color de
  los ojos. Estos genes pueden ser los dos recesivos, los dos dominantes o tener un
  gen recesivo y uno dominante. Llamaremos
  \[
    X = \text{Gen Dominante}, \qquad Y = \text{Gen Recesivo.}
  \]
  Cada uno de los padres transmite a sus descendientes un s\'olo gen. Adem\'as, es
  igualmente probable que se transmita cualquiera de los dos genes que porta el
  progenitor a sus descendientes.

  \smallskip
  \begin{enumerate}[label=\alph*)]
    \item Defina una cadena de Markov donde el estado del sistema sea la cadena
          gen\'etica de los dos individuos de la poblaci\'on. Dibuje el grafo asociado
          y calcule las probabilidades de transici\'on.
    \item Defina las clases, los tipos de estados en cada clase y su correspondiente
          per\'iodo.
    \item ?`Cu\'al es la probabilidad que, en el largo plazo, s\'olo haya individuos
          con genes recesivos si originalmente se parte con dos individuos id\'enticos
          con carga gen\'etica $(X,Y)$?
  \end{enumerate}
\end{enunciado}

\vspace{6pt}

% ═══════════════════════════════════════════════════════════════════════════
%  PARTE a)
% ═══════════════════════════════════════════════════════════════════════════
\parteheader{Parte a) \textperiodcentered\ Espacio de estados, grafo y matriz de transici\'on}

\textbf{Espacio de estados.}\quad
Cada individuo puede tener tres genotipos: $XX$, $XY$ o $YY$. El estado de la
poblaci\'on es el par \emph{no ordenado} de genotipos de los dos individuos; hay
seis estados posibles:

\begin{center}
  \begin{tabular}{cll}
    \toprule
    Estado & Par gen\'etico & Descripci\'on                                     \\
    \midrule
    $E_0$  & $\{XX,\,XX\}$  & Ambos homocigotos dominantes                      \\
    $E_1$  & $\{XX,\,XY\}$  & Uno homocigoto dominante, uno heterocigoto        \\
    $E_2$  & $\{XX,\,YY\}$  & Uno homocigoto dominante, uno homocigoto recesivo \\
    $E_3$  & $\{XY,\,XY\}$  & Ambos heterocigotos                               \\
    $E_4$  & $\{XY,\,YY\}$  & Uno heterocigoto, uno homocigoto recesivo         \\
    $E_5$  & $\{YY,\,YY\}$  & Ambos homocigotos recesivos                       \\
    \bottomrule
  \end{tabular}
\end{center}

\medskip
\textbf{Mecanismo de reproducci\'on.}\quad
Dados dos padres $P_1$ y $P_2$, cada hijo recibe un gen elegido
\emph{uniformemente} de $P_1$ y un gen elegido \emph{uniformemente} de $P_2$.
Los dos hijos se producen de forma \emph{independiente}.

\medskip
\textbf{C\'alculo de probabilidades de transici\'on.}\quad
Se analiza cada estado de origen:

\medskip
\noindent\textit{Desde $E_0 = \{XX,XX\}$:}\quad
Ambos padres tienen genotipo $XX$, por lo que siempre transmiten $X$. Todo hijo
es necesariamente $XX$.
\[
  \pij{E_0}{E_0} = 1.
\]

\noindent\textit{Desde $E_5 = \{YY,YY\}$:}\quad An\'alogo al anterior. Todo hijo es $YY$.
\[
  \pij{E_5}{E_5} = 1.
\]

\noindent\textit{Desde $E_1 = \{XX,XY\}$:}\quad
$P_1 = XX$ siempre env\'ia $X$. $P_2 = XY$ env\'ia $X$ con prob.\ $\tfrac{1}{2}$ e
$Y$ con prob.\ $\tfrac{1}{2}$. El genotipo de cada hijo es:
\[
  XX \text{ con prob.\ } \tfrac{1}{2},\qquad XY \text{ con prob.\ } \tfrac{1}{2}.
\]
Con dos hijos independientes, la distribuci\'on del par resultante es:
\begin{align*}
  \pij{E_1}{E_0} & = \left(\tfrac{1}{2}\right)^2 = \tfrac{1}{4},         \\
  \pij{E_1}{E_1} & = 2\cdot\tfrac{1}{2}\cdot\tfrac{1}{2} = \tfrac{1}{2}, \\
  \pij{E_1}{E_3} & = \left(\tfrac{1}{2}\right)^2 = \tfrac{1}{4}.
\end{align*}

\noindent\textit{Desde $E_4 = \{XY,YY\}$:}\quad
$P_1 = XY$ env\'ia $X$ o $Y$ con prob.\ $\tfrac{1}{2}$ cada uno. $P_2 = YY$ siempre
env\'ia $Y$. Cada hijo es $XY$ con prob.\ $\tfrac{1}{2}$ o $YY$ con prob.\ $\tfrac{1}{2}$.
Par resultado:
\[
  \pij{E_4}{E_3} = \tfrac{1}{4},\qquad \pij{E_4}{E_4} = \tfrac{1}{2},\qquad
  \pij{E_4}{E_5} = \tfrac{1}{4}.
\]

\noindent\textit{Desde $E_2 = \{XX,YY\}$:}\quad
$P_1 = XX$ siempre env\'ia $X$; $P_2 = YY$ siempre env\'ia $Y$. Todo hijo es
$XY$ con certeza.
\[
  \pij{E_2}{E_3} = 1.
\]

\noindent\textit{Desde $E_3 = \{XY,XY\}$:}\quad
Cada padre $XY$ env\'ia $X$ o $Y$ con prob.\ $\tfrac{1}{2}$. El genotipo de cada
hijo sigue:
\[
  \Prob{\text{hijo}=XX} = \tfrac{1}{4},\quad
  \Prob{\text{hijo}=XY} = \tfrac{1}{2},\quad
  \Prob{\text{hijo}=YY} = \tfrac{1}{4}.
\]
Para el par de dos hijos independientes:
\begin{align*}
  \pij{E_3}{E_0} & = \left(\tfrac{1}{4}\right)^2 = \tfrac{1}{16},         \\[2pt]
  \pij{E_3}{E_1} & = 2\cdot\tfrac{1}{4}\cdot\tfrac{1}{2} = \tfrac{4}{16}, \\[2pt]
  \pij{E_3}{E_2} & = 2\cdot\tfrac{1}{4}\cdot\tfrac{1}{4} = \tfrac{2}{16}, \\[2pt]
  \pij{E_3}{E_3} & = \left(\tfrac{1}{2}\right)^2 = \tfrac{4}{16},         \\[2pt]
  \pij{E_3}{E_4} & = 2\cdot\tfrac{1}{2}\cdot\tfrac{1}{4} = \tfrac{4}{16}, \\[2pt]
  \pij{E_3}{E_5} & = \left(\tfrac{1}{4}\right)^2 = \tfrac{1}{16}.
\end{align*}

\medskip
\textbf{Matriz de transici\'on.}\quad
Ordenando los estados como $E_0,E_1,E_2,E_3,E_4,E_5$:

\[
  \Pmat =
  \begin{pmatrix}
    1             & 0             & 0             & 0             & 0             & 0             \\[4pt]
    \tfrac{1}{4}  & \tfrac{1}{2}  & 0             & \tfrac{1}{4}  & 0             & 0             \\[4pt]
    0             & 0             & 0             & 1             & 0             & 0             \\[4pt]
    \tfrac{1}{16} & \tfrac{4}{16} & \tfrac{2}{16} & \tfrac{4}{16} & \tfrac{4}{16} & \tfrac{1}{16} \\[4pt]
    0             & 0             & 0             & \tfrac{1}{4}  & \tfrac{1}{2}  & \tfrac{1}{4}  \\[4pt]
    0             & 0             & 0             & 0             & 0             & 1
  \end{pmatrix}
\]

\medskip
\textbf{Grafo de transici\'on.}

\begin{center}
  \begin{tikzpicture}[
      every node/.style={circle, draw=sectionbg, fill=sectionbg!10,
          minimum size=1.0cm, font=\small\sffamily},
      every edge/.style={draw, -Stealth, thick},
      lbl/.style={draw=none, fill=none, font=\footnotesize, inner sep=2pt},
    ]

    % Dos filas: absorbentes a la izquierda, transitoria al centro y derecha
    %   Fila superior: E0 — E1 — E2
    %   Fila inferior: E5 — E4 — E3
    \node (E0) at (0,   2.6) {$E_0$};
    \node (E1) at (3.6, 2.6) {$E_1$};
    \node (E2) at (7.2, 3.9) {$E_2$};
    \node (E5) at (0,   0)   {$E_5$};
    \node (E4) at (3.6, 0)   {$E_4$};
    \node (E3) at (7.2, 1.3)   {$E_3$};

    % ── Auto-lazos ──────────────────────────────────────────────────────
    \draw (E0) edge[loop left,  looseness=5] node[lbl, left]  {$1$}             (E0);
    \draw (E5) edge[loop left,  looseness=5] node[lbl, left]  {$1$}             (E5);
    \draw (E1) edge[loop above, looseness=5] node[lbl, above] {$\tfrac{1}{2}$}  (E1);
    \draw (E3) edge[loop right, looseness=5] node[lbl, right] {$\tfrac{4}{16}$} (E3);
    \draw (E4) edge[loop below, looseness=5] node[lbl, below] {$\tfrac{1}{2}$}  (E4);

    % ── Fila superior: E1 → E0 ──────────────────────────────────────────
    \draw (E1) edge
      node[lbl, above] {$\tfrac{1}{4}$} (E0);

    % ── Columna derecha: E2 ↔ E3 (bend left separa los dos sentidos) ───
    \draw (E2) edge[bend left=12]
      node[lbl, right=3pt] {$1$}             (E3);
    \draw (E3) edge[bend left=12]
      node[lbl, left=3pt]  {$\tfrac{2}{16}$} (E2);

    % ── Diagonal: E3 ↔ E1 (bend left separa los dos sentidos) ──────────
    \draw (E3) edge[bend left=12]
      node[lbl, above, pos=0.42] {$\tfrac{4}{16}$} (E1);
    \draw (E1) edge[bend left=12]
      node[lbl, below, pos=0.42] {$\tfrac{1}{4}$}  (E3);

    % ── Larga diagonal: E3 → E0 (arco por encima del diagrama) ─────────
    \draw (E3) edge[out=128, in=18]
      node[lbl, above, pos=0.42] {$\tfrac{1}{16}$} (E0);

    % ── Fila inferior: E3 ↔ E4 (bend left separa los dos sentidos) ─────
    \draw (E3) edge[bend left=12]
      node[lbl, above] {$\tfrac{4}{16}$} (E4);
    \draw (E4) edge[bend left=12]
      node[lbl, below] {$\tfrac{1}{4}$}  (E3);

    % ── Fila inferior: E4 → E5 ──────────────────────────────────────────
    \draw (E4) edge
      node[lbl, below] {$\tfrac{1}{4}$} (E5);

    % ── Larga inferior: E3 → E5 (arco por debajo del diagrama) ─────────
    \draw (E3) edge[out=207, in=-27]
      node[lbl, below, pos=0.42] {$\tfrac{1}{16}$} (E5);

  \end{tikzpicture}
\end{center}

% ═══════════════════════════════════════════════════════════════════════════
%  PARTE b)
% ═══════════════════════════════════════════════════════════════════════════
\parteheader{Parte b) \textperiodcentered\ Clases de comunicaci\'on, tipos y per\'iodos}

\textbf{Comunicaci\'on entre estados.}\quad
Se dice que $i \leftrightarrow j$ si $i \to j$ y $j \to i$ (ambos alcanzables entre
s\'i). Se verifican las siguientes rutas:

\begin{itemize}
  \item $E_1 \leftrightarrow E_3$: directo en un paso ($\pij{E_1}{E_3}=\tfrac{1}{4}>0$
        y $\pij{E_3}{E_1}=\tfrac{4}{16}>0$).
  \item $E_2 \leftrightarrow E_3$: $E_2\to E_3$ en un paso; $E_3\to E_2$ en un paso.
  \item $E_4 \leftrightarrow E_3$: directo en un paso.
  \item Por transitividad: $E_1\leftrightarrow E_2\leftrightarrow E_3\leftrightarrow E_4$.
\end{itemize}

$E_0$ y $E_5$ son \emph{absorbentes}: una vez alcanzados no se puede salir, y por
tanto no comunican con ning\'un otro estado.

\medskip
\textbf{Clases de comunicaci\'on:}

\begin{center}
  \begin{tabular}{clll}
    \toprule
    Clase           & Estados                  & Tipo                    & Per\'iodo \\
    \midrule
    $\mathcal{C}_0$ & $\{E_0\}$                & Recurrente (absorbente) & $d = 1$   \\
    $\mathcal{C}_1$ & $\{E_1, E_2, E_3, E_4\}$ & Transitorio             & $d = 1$   \\
    $\mathcal{C}_2$ & $\{E_5\}$                & Recurrente (absorbente) & $d = 1$   \\
    \bottomrule
  \end{tabular}
\end{center}

\textbf{Justificaci\'on del tipo transitorio de $\mathcal{C}_1$:}\quad
Todo estado en $\{E_1,E_2,E_3,E_4\}$ puede alcanzar $E_0$ o $E_5$ con probabilidad
positiva (por ejemplo, $E_3\to E_0$ con prob.\ $\tfrac{1}{16}$). Una vez absorbidos
en $\mathcal{C}_0$ o $\mathcal{C}_2$, no es posible regresar. Por tanto, la
probabilidad de retornar a cualquier estado de $\mathcal{C}_1$ es estrictamente
menor a $1$: todos los estados son \emph{transitorios}.

\medskip
\textbf{Per\'iodo de los estados en $\mathcal{C}_1$:}\quad
El per\'iodo de un estado $i$ es $d(i)=\gcd\{n\ge 1: \pij{i}{i}^{(n)}>0\}$,
y es el mismo para todos los estados de una clase.

\begin{itemize}
  \item $E_1$, $E_3$, $E_4$ tienen auto-lazo ($\pij{i}{i}>0$), luego $n=1$ es
        alcanzable. Esto implica $d=1$ (aperiódico).
  \item $E_2$ no tiene auto-lazo, pero puede retornar en 2 pasos
        ($E_2\to E_3\to E_2$) y en 3 pasos ($E_2\to E_3\to E_3\to E_2$, usando
        el auto-lazo de $E_3$). Como $\gcd(2,3)=1$, se concluye $d(E_2)=1$.
\end{itemize}

Todos los estados de $\mathcal{C}_1$ son \textbf{aperiódicos} ($d=1$).

% ═══════════════════════════════════════════════════════════════════════════
%  PARTE c)
% ═══════════════════════════════════════════════════════════════════════════
\parteheader{Parte c) \textperiodcentered\ Probabilidad de absorci\'on en $E_5$ desde $E_3$}

El estado inicial es $E_3 = \{XY,XY\}$ (dos individuos id\'enticos con genotipo
$XY$). Se quiere calcular la probabilidad de que la cadena sea absorbida en $E_5$
(``todos recesivos'') a largo plazo.

\medskip
\textbf{Probabilidades de absorci\'on $f_i$.}\quad
Dado que la cadena puede pasar por varios estados transitorios antes de absorberse,
no es posible calcular $f_3$ directamente desde la definici\'on. La estrategia es
definir, para \emph{cada} estado $E_i$,
\[
  f_i \;=\; \Prob{\text{absorci\'on en }E_5 \mid X_0 = E_i},
\]
es decir, $f_i$ es la probabilidad de que la cadena, \textbf{partiendo del estado
  $E_i$}, acabe atrapada en $E_5$ en alg\'un tiempo futuro (potencialmente muy lejano).
Al conocer $f_i$ para todos los estados, se puede recuperar el valor buscado
$f_3$ resolviendo un sistema de ecuaciones lineales.

\medskip
\textbf{Planteamiento del sistema.}\quad
Las condiciones de frontera son inmediatas: si ya se est\'a en un estado
absorbente, la probabilidad de llegar a $E_5$ es conocida:
\[
  f_0 = 0 \quad (\text{absorbido en }E_0\text{, nunca llega a }E_5), \qquad
  f_5 = 1 \quad (\text{ya est\'a en }E_5).
\]
Para los estados transitorios, por la \emph{ley de probabilidad total} (condicionando
en el primer paso):
\begin{align}
  f_1 & = \pij{E_1}{E_0}\,f_0 + \pij{E_1}{E_1}\,f_1 + \pij{E_1}{E_3}\,f_3
  = \tfrac{1}{4}\cdot 0 + \tfrac{1}{2}\,f_1 + \tfrac{1}{4}\,f_3,          \\[4pt]
  f_2 & = \pij{E_2}{E_3}\,f_3 = f_3,                                      \\[4pt]
  f_4 & = \pij{E_4}{E_3}\,f_3 + \pij{E_4}{E_4}\,f_4 + \pij{E_4}{E_5}\,f_5
  = \tfrac{1}{4}\,f_3 + \tfrac{1}{2}\,f_4 + \tfrac{1}{4},                 \\[4pt]
  f_3 & = \tfrac{1}{16}\,f_0 + \tfrac{4}{16}\,f_1 + \tfrac{2}{16}\,f_2
  + \tfrac{4}{16}\,f_3 + \tfrac{4}{16}\,f_4 + \tfrac{1}{16}\,f_5.
\end{align}

\medskip
\textbf{Resoluci\'on paso a paso.}

\noindent\textit{De la ecuaci\'on (1):}
\[
  f_1 - \tfrac{1}{2}\,f_1 = \tfrac{1}{4}\,f_3 \implies
  \tfrac{1}{2}\,f_1 = \tfrac{1}{4}\,f_3 \implies
  \boxed{f_1 = \tfrac{1}{2}\,f_3.}
\]

\noindent\textit{De la ecuaci\'on (2):} $\boxed{f_2 = f_3.}$

\noindent\textit{De la ecuaci\'on (3):}
\[
  f_4 - \tfrac{1}{2}\,f_4 = \tfrac{1}{4}\,f_3 + \tfrac{1}{4} \implies
  \tfrac{1}{2}\,f_4 = \tfrac{1}{4}\,f_3 + \tfrac{1}{4} \implies
  \boxed{f_4 = \tfrac{1}{2}\,f_3 + \tfrac{1}{2}.}
\]

\noindent\textit{Sustituyendo en la ecuaci\'on (4)}, con $f_0=0$ y $f_5=1$:
\begin{align*}
  f_3 & = 0 + \tfrac{4}{16}\cdot\tfrac{1}{2}\,f_3
  + \tfrac{2}{16}\cdot f_3
  + \tfrac{4}{16}\cdot f_3
  + \tfrac{4}{16}\!\left(\tfrac{1}{2}\,f_3+\tfrac{1}{2}\right)
  + \tfrac{1}{16}                                                      \\[4pt]
      & = \tfrac{2}{16}\,f_3 + \tfrac{2}{16}\,f_3 + \tfrac{4}{16}\,f_3
  + \tfrac{2}{16}\,f_3 + \tfrac{2}{16} + \tfrac{1}{16}                 \\[4pt]
      & = \tfrac{10}{16}\,f_3 + \tfrac{3}{16}.
\end{align*}
\[
  f_3 - \tfrac{10}{16}\,f_3 = \tfrac{3}{16} \implies
  \tfrac{6}{16}\,f_3 = \tfrac{3}{16} \implies
  \boxed{f_3 = \tfrac{1}{2}.}
\]

\medskip
\textbf{Respuesta.}\quad La probabilidad de que en el largo plazo solo haya
individuos con genes recesivos, partiendo de $\{XY,XY\}$, es
\[
  \Prob{\text{absorci\'on en }E_5 \mid X_0 = E_3} = \dfrac{1}{2}.
\]

\begin{nota}
  \textbf{Interpretaci\'on por simetr\'ia.} El modelo trata a los genes $X$ e $Y$ de
  forma completamente sim\'etrica. En el estado $E_3 = \{XY,XY\}$ hay exactamente
  2 genes dominantes y 2 recesivos sobre 4 totales (frecuencia recesiva $= 50\%$).
  Por simetr\'ia, la probabilidad de absorci\'on en ``todos dominantes'' es tambi\'en
  $\tfrac{1}{2}$, y ambas suman 1 (la absorci\'on es cierta, pues los estados
  transitorios se visitan un n\'umero finito de veces). En general, la probabilidad
  de absorci\'on en $E_5$ desde cualquier estado transitorio es igual a la
  \emph{frecuencia inicial de genes recesivos}: $f_1=\tfrac{1}{4}$,
  $f_2=f_3=\tfrac{1}{2}$, $f_4=\tfrac{3}{4}$.
\end{nota}

\vspace{10pt}
{\color{subsecbg!60}\hrule}
\vspace{6pt}

% ═══════════════════════════════════════════════════════════════════════
\begin{enunciado}{Problema 7 \textperiodcentered\ Centro hospitalario: instalaci\'on de equipos}
  En un peque\~no centro hospitalario se tiene la urgencia de instalar equipos nuevos.
  Estos equipos son muy costosos y se deben manejar con mucho cuidado, por lo que se
  necesita que el establecimiento est\'e vac\'io al momento de la instalaci\'on.
  Actualmente hay $M$ pacientes en el centro y, con el fin de poder proceder a la
  instalaci\'on, no se recibir\'an m\'as pacientes hasta que \'esta se realice.

  Cada ma\~nana un doctor eval\'ua la condici\'on de los pacientes. Se ha determinado
  que cada paciente tiene una probabilidad $p$ de estar rehabilitado y salir del
  centro, y una probabilidad $(1-p)$ de seguir internado, \emph{independientemente}
  de lo que ocurra con los dem\'as pacientes. Nadie puede ingresar al centro hasta
  despu\'es de instalados los equipos.

  \smallskip
  \begin{enumerate}[label=\alph*)]
    \item Muestre que el sistema descrito puede ser modelado como una cadena de
          Markov en tiempo discreto, dibuje el grafo correspondiente, identifique
          las clases y clasifique sus estados.\\[2pt]
          \textit{Hint:} Calcule la probabilidad
          $\Pr\!\bigl[X(n)=r\mid X(n-1)=k\bigr]$, con $X(n)$ representando la
          cantidad de personas en el centro en la semana $n$.
    \item Si el sistema tiene inicialmente $M$ pacientes, ?`cu\'al es la probabilidad
          de que alg\'un d\'ia haya $M-1$ pacientes? ?`Cu\'al es la probabilidad de
          que alg\'un d\'ia se puedan instalar los equipos? Encuentre estas
          probabilidades y fund\'amente adecuadamente sus respuestas.
  \end{enumerate}
\end{enunciado}

\vspace{6pt}

% ═══════════════════════════════════════════════════════════════════════════
%  PARTE a)
% ═══════════════════════════════════════════════════════════════════════════
\parteheader{Parte a) \textperiodcentered\ Modelo como cadena de Markov, grafo y clasificaci\'on}

\textbf{Definici\'on del proceso.}\quad
Sea $X(n)$ el n\'umero de pacientes en el centro al inicio del d\'ia $n$. El espacio
de estados es $\mathcal{S} = \{0, 1, 2, \ldots, M\}$, y no se admiten nuevos
pacientes, por lo que el proceso s\'olo puede decrecer.

\medskip
\textbf{C\'alculo de probabilidades de transici\'on.}\quad
Dado $X(n-1)=k$, cada uno de los $k$ pacientes permanece con probabilidad $(1-p)$
o es dado de alta con probabilidad $p$, de forma \emph{independiente}. El n\'umero
de pacientes que queda al d\'ia siguiente es por tanto
\[
  X(n)\;\big|\;X(n-1)=k \;\sim\; \mathrm{Binomial}\!\bigl(k,\;1-p\bigr).
\]
La probabilidad de transici\'on de $k$ a $r$ ($0\le r\le k$) es:
\[
  \pij{kr} \;=\; \Pr\!\bigl[X(n)=r\mid X(n-1)=k\bigr]
  \;=\; \binom{k}{r}(1-p)^{r}\,p^{\,k-r},
\]
y $\pij{kr}=0$ para $r>k$ (los pacientes s\'olo pueden marcharse).

\medskip
\textbf{Propiedad de Markov.}\quad
El estado futuro $X(n)$ depende del pasado \emph{\'unicamente} a trav\'es del estado
actual $X(n-1)=k$: dada la cantidad actual de pacientes, el mecanismo de alta de
d\'ias anteriores no aporta informaci\'on adicional. Las probabilidades de
transici\'on son adem\'as hom\'nogeneas en el tiempo. Por tanto, el proceso es una
\textbf{cadena de Markov homog\'enea en tiempo discreto}.

\medskip
\textbf{Grafo de transici\'on} (esquema para $M$ general):

\begin{center}
  \begin{tikzpicture}[
      every node/.style={circle, draw=sectionbg, fill=sectionbg!10,
          minimum size=1.0cm, font=\small\sffamily},
      every edge/.style={draw, -Stealth, thick},
      lbl/.style={draw=none, fill=none, font=\scriptsize, inner sep=2pt},
    ]

    % Estados explícitos
    \node (s0) at (0,   0) {$0$};
    \node (s1) at (2.4, 0) {$1$};
    \node[draw=none, fill=none, minimum size=0pt] (sd1) at (4.4, 0) {\large$\cdots$};
    \node (sk) at (6.4, 0) {$k$};
    \node[draw=none, fill=none, minimum size=0pt] (sd2) at (8.4, 0) {\large$\cdots$};
    \node (sM) at (10.4, 0) {$M$};

    % Auto-lazos
    \draw (s0) edge[loop above, looseness=5] node[lbl, above] {$1$}         (s0);
    \draw (s1) edge[loop above, looseness=5] node[lbl, above] {$(1-p)$}     (s1);
    \draw (sk) edge[loop above, looseness=5] node[lbl, above] {$(1-p)^k$}   (sk);
    \draw (sM) edge[loop above, looseness=5] node[lbl, above] {$(1-p)^M$}   (sM);

    % 1 → 0
    \draw (s1) edge
      node[lbl, above] {$p$} (s0);

    % k → 1 (transición representativa intermedia)
    \draw (sk) edge[bend left=22]
      node[lbl, above, pos=0.5] {$\binom{k}{1}(1-p)\,p^{k-1}$} (s1);

    % k → 0
    \draw (sk) edge[bend right=25]
      node[lbl, below, pos=0.5] {$p^k$} (s0);

    % M → k (transición representativa)
    \draw (sM) edge[bend left=20]
      node[lbl, above, pos=0.45] {$\binom{M}{k}(1-p)^k p^{M-k}$} (sk);

    % M → 0
    \draw (sM) edge[out=215, in=-35]
      node[lbl, below, pos=0.42] {$p^M$} (s0);

  \end{tikzpicture}
\end{center}

\medskip
\textbf{Clases de comunicaci\'on y clasificaci\'on de estados.}\quad
Dado que el proceso s\'olo puede \emph{decrecer} (nunca entran nuevos pacientes),
ning\'un estado $k\ge 1$ puede ser alcanzado desde ning\'un estado $j<k$.

\begin{itemize}
  \item \textbf{Estado 0:} una vez vac\'io el centro, permanece vac\'io
        ($\pij{00}=1$). Es un \emph{estado absorbente}: clase recurrente
        $\mathcal{C}_0 = \{0\}$, periodo $d=1$.
  \item \textbf{Estados $k=1,\ldots,M$:} desde $k$ se puede alcanzar cualquier
        $j<k$, pero desde $j$ no se puede retornar a $k$. Por tanto $k$ no
        comunica con ning\'un otro estado y forma su propio singleton
        $\{k\}$, que es \emph{transitorio} y \emph{aperi\'odico} ($d=1$,
        pues $\pij{kk}=(1-p)^k>0$).
\end{itemize}

\begin{center}
  \begin{tabular}{cll}
    \toprule
    Clase           & Estados                      & Tipo                           \\
    \midrule
    $\mathcal{C}_0$ & $\{0\}$                      & Recurrente (absorbente), $d=1$ \\
    $\mathcal{C}_k$ & $\{k\}$, para $k=1,\ldots,M$ & Transitorio, $d=1$             \\
    \bottomrule
  \end{tabular}
\end{center}

% ═══════════════════════════════════════════════════════════════════════════
%  PARTE b)
% ═══════════════════════════════════════════════════════════════════════════
\parteheader{Parte b) \textperiodcentered\ Probabilidades de alcanzar $M{-}1$ y de instalar los equipos}

\textbf{Probabilidad de que alg\'un d\'ia haya $M-1$ pacientes.}\quad
Sea
\[
  f \;=\; \Prob{\exists\,n\ge 1:\,X(n)=M-1 \mid X(0)=M}.
\]
Dado que el proceso s\'olo decrece, si alguna vez cae por debajo de $M-1$ nunca
podr\'a volver a $M-1$. Condicionando en el primer paso:
\begin{align*}
  f & = \Prob{X(1)=M-1\mid X(0)=M}
  + \Prob{X(1)=M\mid X(0)=M}\cdot f
  + \underbrace{\sum_{r=0}^{M-2}\Prob{X(1)=r\mid X(0)=M}\cdot 0}_{\text{cae debajo de }M-1:\text{ imposible regresar}}.
\end{align*}
Definiendo
\[
  \alpha \;=\; \Prob{X(1)=M-1\mid X(0)=M}
  \;=\; \binom{M}{M-1}(1-p)^{M-1}p
  \;=\; M\,p\,(1-p)^{M-1},
\]
\[
  \beta  \;=\; \Prob{X(1)=M\mid X(0)=M}
  \;=\; (1-p)^M,
\]
la ecuaci\'on se reduce a $f = \alpha + \beta f$, de donde:
\[
  \boxed{f \;=\; \frac{\alpha}{1-\beta}
    \;=\; \frac{M\,p\,(1-p)^{M-1}}{1-(1-p)^M}.}
\]

\begin{nota}
  Esta f\'ormula se puede interpretar como la probabilidad de dar exactamente un paso
  de $M$ a $M-1$, condicionada a que el sistema \emph{s\'i} abandone el estado $M$
  (evento $1-\beta$). Para $M=1$ se obtiene $f=1$ (trivialmente, el \'unico paciente
  debe pasar por 0, que requiere estar en 1 antes).
\end{nota}

\medskip
\textbf{Probabilidad de instalar los equipos.}\quad
Se busca $\Prob{\exists\,n\ge 1:\,X(n)=0 \mid X(0)=M}$. Se mostrar\'a que esta
probabilidad es \textbf{igual a 1} para cualquier $p\in(0,1)$.

\medskip
\textit{Argumento por esperanza.}\quad
Cada uno de los $M$ pacientes presentes al inicio permanece en el centro al d\'ia
$n$ con probabilidad $(1-p)^n$, independientemente de los dem\'as. Por linealidad
de la esperanza:
\[
  \E{X(n)\mid X(0)=M} \;=\; M\,(1-p)^n \;\xrightarrow{n\to\infty}\; 0,
  \qquad \text{para } p\in(0,1).
\]
Como $X(n)\ge 0$ es entero, la desigualdad de Markov da:
\[
  \Prob{X(n)\ge 1 \mid X(0)=M}
  \;\le\; \E{X(n)\mid X(0)=M}
  \;=\; M\,(1-p)^n \;\to\; 0.
\]
Por tanto $\Prob{X(n)=0\mid X(0)=M}\to 1$, lo que implica que la absorci\'on en
el estado 0 ocurre con probabilidad 1:
\[
  \boxed{\Prob{\text{equipos instalados}\mid X(0)=M}
    \;=\; \Prob{\exists\,n:\,X(n)=0 \mid X(0)=M} \;=\; 1.}
\]

\begin{nota}
  El resultado se puede verificar tambi\'en notando que desde cualquier estado
  $k\ge 1$ existe una probabilidad $p^k>0$ de llegar a 0 en un solo paso
  (todos los pacientes salen el mismo d\'ia). Como esta probabilidad est\'a
  uniformemente acotada por $p^M>0$, la probabilidad de \emph{no} absorci\'on
  en $n$ pasos est\'a acotada por $(1-p^M)^n\to 0$.
\end{nota}

\vspace{10pt}
{\color{subsecbg!60}\hrule}
\vspace{6pt}

% ═══════════════════════════════════════════════════════════════════════
\begin{enunciado}{P35 \textperiodcentered\ Empresa de arriendo de autos: Santiago--Vi\~na del Mar}
  Considere una empresa automovilística que arrienda autos entre Santiago y Vi\~na del
  Mar. La empresa dispone de $N=3$ autos en total. Las funciones de probabilidad de
  la demanda diaria son independientes entre ciudades y a lo largo del tiempo:

  \smallskip
  \begin{center}
    \begin{tabular}{lcccccc}
      \toprule
      Cantidad demandada & 0    & 1    & 2    & 3    & 4    & 5    \\
      \midrule
      Santiago           & 0.10 & 0.20 & 0.50 & 0.10 & 0.05 & 0.05 \\
      Vi\~na del Mar     & 0.30 & 0.50 & 0.10 & 0.10 & 0    & 0    \\
      \bottomrule
    \end{tabular}
  \end{center}

  \smallskip
  Los autos arrendados viajan a la otra ciudad ese mismo d\'ia (disponibles allí al
  final del d\'ia). La demanda no satisfecha se pierde. Sea $X_n$ el n\'umero de autos
  en Santiago al comienzo del d\'ia $n$.

  \smallskip
  \begin{enumerate}[label=\alph*)]
    \item Obtenga la matriz de probabilidades de transici\'on en una etapa.
    \item Suponga que $X_1=2$. Obtenga el valor esperado de la demanda insatisfecha
          por la empresa durante el primer d\'ia.
    \item La empresa decide implementar la siguiente pol\'itica nocturna: si al final
          del d\'ia cualquiera de las dos ciudades se queda sin autos, se traslada un
          auto desde la otra ciudad.
          \begin{enumerate}[label=\roman*.]
            \item Obtenga la nueva matriz de transici\'on en una etapa.
            \item Suponga que ha obtenido el vector estacionario $\pi$ de este proceso.
                  Encuentre una expresi\'on para el n\'umero promedio de traslados nocturnos
                  efectuados en el largo plazo.
          \end{enumerate}
  \end{enumerate}
\end{enunciado}

\vspace{6pt}

% ───────────────────────────────────────────────────────────────────────
\parteheader{Parte a) \textperiodcentered\ Matriz de transici\'on en una etapa}

\textbf{Din\'amica del sistema.}\quad
Si al inicio del d\'ia $n$ hay $i$ autos en Santiago (y $3-i$ en Vi\~na), la cantidad
satisfecha en cada ciudad queda limitada por el stock disponible:
\[
  A = \min(D_S,\,i), \qquad B = \min(D_V,\,3-i).
\]
Los $A$ autos arrendados en Santiago viajan a Vi\~na ese d\'ia; los $B$ arrendados en
Vi\~na viajan a Santiago. Por tanto:
\[
  X_{n+1} = i - A + B.
\]
Como $D_S \perp D_V$ (independientes), las distribuciones de $A$ y $B$ se obtienen
directamente de las tablas de demanda.

\medskip
\noindent\textit{Estado $i=0$ (0 autos en Santiago):}\quad
$A=0$ con certeza. $B = \min(D_V,3)$, luego $X_{n+1}=B$.
\[
  p_{00}=0.30,\quad p_{01}=0.50,\quad p_{02}=0.10,\quad
  p_{03}=\Prob{D_V\ge 3}=0.10.
\]

\noindent\textit{Estado $i=3$ (3 autos en Santiago):}\quad
$B=0$ con certeza. $A=\min(D_S,3)$, luego $X_{n+1}=3-A$.
\[
  p_{30}=\Prob{D_S\ge 3}=0.20,\quad p_{31}=0.50,\quad p_{32}=0.20,\quad p_{33}=0.10.
\]

\noindent\textit{Estado $i=1$ (1 auto en Santiago, 2 en Vi\~na):}\quad
$A\in\{0,1\}$, $B\in\{0,1,2\}$. Tabla de probabilidades conjuntas de $X_{n+1}=1-A+B$:

\begin{center}
  \begin{tabular}{l|ccc|l}
    \toprule
                                 & $B=0$                                                         & $B=1$       & $B=2$       & \\
                                 & (0.30)                                                        & (0.50)      & (0.20)      & \\
    \midrule
    $A=0$ (0.10)                 & $X=1$, 0.03                                                   & $X=2$, 0.05 & $X=3$, 0.02 & \\
    $A=1$ (0.90)                 & $X=0$, 0.27                                                   & $X=1$, 0.45 & $X=2$, 0.18 & \\
    \midrule
    \multicolumn{4}{l}{Sumando:} & $p_{10}=0.27$,\ $p_{11}=0.48$,\ $p_{12}=0.23$,\ $p_{13}=0.02$                               \\
    \bottomrule
  \end{tabular}
\end{center}

\noindent\textit{Estado $i=2$ (2 autos en Santiago, 1 en Vi\~na):}\quad
$A\in\{0,1,2\}$, $B\in\{0,1\}$. Tabla de $X_{n+1}=2-A+B$:

\begin{center}
  \begin{tabular}{l|cc|l}
    \toprule
                                 & $B=0$ (0.30)                                                  & $B=1$ (0.70) & \\
    \midrule
    $A=0$ (0.10)                 & $X=2$, 0.03                                                   & $X=3$, 0.07  & \\
    $A=1$ (0.20)                 & $X=1$, 0.06                                                   & $X=2$, 0.14  & \\
    $A=2$ (0.70)                 & $X=0$, 0.21                                                   & $X=1$, 0.49  & \\
    \midrule
    \multicolumn{3}{l}{Sumando:} & $p_{20}=0.21$,\ $p_{21}=0.55$,\ $p_{22}=0.17$,\ $p_{23}=0.07$                  \\
    \bottomrule
  \end{tabular}
\end{center}

\medskip
\textbf{Matriz de transici\'on} (filas = estado origen, columnas = estado destino):
\[
  \mathbf{P} =
  \bordermatrix{
    & 0    & 1    & 2    & 3    \cr
    i=0    & 0.30 & 0.50 & 0.10 & 0.10 \cr
    i=1    & 0.27 & 0.48 & 0.23 & 0.02 \cr
    i=2    & 0.21 & 0.55 & 0.17 & 0.07 \cr
    i=3    & 0.20 & 0.50 & 0.20 & 0.10 \cr
  }
\]

\medskip
\textbf{Clasificaci\'on.}\quad
Todos los elementos de $\mathbf{P}$ son estrictamente positivos, por lo que desde
cualquier estado se puede alcanzar cualquier otro en un solo paso: la cadena es
\emph{irreducible}. Adem\'as, todos los estados tienen auto-lazo ($p_{ii}>0$), por
lo que son \emph{aperi\'odicos} ($d=1$). Dado que el espacio de estados es finito,
todos los estados son \textbf{recurrentes positivos} y existe una \'unica
distribuci\'on estacionaria $\pi$.

% ───────────────────────────────────────────────────────────────────────
\parteheader{Parte b) \textperiodcentered\ Demanda insatisfecha esperada con $X_1=2$}

Con $X_1=2$: hay 2 autos en Santiago y 1 en Vi\~na. La demanda insatisfecha total
en el d\'ia 1 es:
\[
  U = (D_S - 2)^+ + (D_V - 1)^+,
  \qquad \text{donde } (x)^+ = \max(x,0).
\]
Dado que $D_S \perp D_V$, $\E{U} = \E{(D_S-2)^+} + \E{(D_V-1)^+}$.

\[
  \E{(D_S-2)^+} = 1\cdot\Prob{D_S=3} + 2\cdot\Prob{D_S=4} + 3\cdot\Prob{D_S=5}
  = 1(0.10)+2(0.05)+3(0.05) = 0.35.
\]
\[
  \E{(D_V-1)^+} = 1\cdot\Prob{D_V=2} + 2\cdot\Prob{D_V=3}
  = 1(0.10)+2(0.10) = 0.30.
\]
\[
  \boxed{\E{U \mid X_1=2} = 0.35 + 0.30 = 0.65 \text{ autos/d\'ia.}}
\]

% ───────────────────────────────────────────────────────────────────────
\parteheader{Parte c) \textperiodcentered\ Pol\'itica de traslado nocturno}

\textbf{i) Nueva matriz de transici\'on.}\quad
La pol\'itica transforma el estado \emph{al final del d\'ia} antes de comenzar el
siguiente:
\[
  Y = X_{n+1}^{\text{raw}} \mapsto X_{n+1} =
  \begin{cases}
    1 & \text{si } Y=0 \;\text{(Vi\~na tiene 3, traslado a Santiago)}, \\
    Y & \text{si } Y=1 \text{ \'o } 2,                                 \\
    2 & \text{si } Y=3 \;\text{(Santiago tiene 3, traslado a Vi\~na)}.
  \end{cases}
\]
Los estados 0 y 3 nunca son estado inicial al d\'ia siguiente: el espacio de
estados efectivo se reduce a $\{1,2\}$. La nueva matriz $\mathbf{P}'$ se obtiene
agrupando columnas de $\mathbf{P}$:
\begin{align*}
  p'_{11} & = p_{10}+p_{11} = 0.27+0.48 = 0.75, &
  p'_{12} & = p_{12}+p_{13} = 0.23+0.02 = 0.25,   \\
  p'_{21} & = p_{20}+p_{21} = 0.21+0.55 = 0.76, &
  p'_{22} & = p_{22}+p_{23} = 0.17+0.07 = 0.24.
\end{align*}
\[
  \mathbf{P}' = \begin{pmatrix} 0.75 & 0.25 \\ 0.76 & 0.24 \end{pmatrix}.
\]

\textbf{ii) Promedio de traslados en el largo plazo.}\quad
Un traslado ocurre cuando el estado \emph{bruto} al final del d\'ia es $Y=0$
(traslado de Vi\~na a Santiago) o $Y=3$ (traslado de Santiago a Vi\~na). En el largo
plazo, la frecuencia de d\'ias con traslado es:
\[
  \bar{T} = \pi_1\bigl(p_{10}+p_{13}\bigr) + \pi_2\bigl(p_{20}+p_{23}\bigr)
  = 0.29\,\pi_1 + 0.28\,\pi_2,
\]
donde $\pi=(\pi_1,\pi_2)$ es la distribuci\'on estacionaria de $\mathbf{P}'$.

\begin{nota}
  Resolviendo $\pi\mathbf{P}'=\pi$, $\pi_1+\pi_2=1$: de $0.25\,\pi_1=0.76\,\pi_2$
  se obtiene $\pi_1=76/101$, $\pi_2=25/101$. Sustituyendo:
  $\bar{T} = 0.29\times(76/101)+0.28\times(25/101) = 2904/10100 \approx 0.288$
  traslados por d\'ia.
\end{nota}

\vspace{10pt}
{\color{subsecbg!60}\hrule}
\vspace{6pt}

% ═══════════════════════════════════════════════════════════════════════
\begin{enunciado}{Problema 2 \textperiodcentered\ Ruina del jugador}
  Un jugador apuesta sucesivas veces en el mismo juego. En cada jugada existe una
  probabilidad $p$ de ganar una unidad y una probabilidad $1-p$ de perder una unidad.
  Las jugadas son independientes. El jugador comienza con una fortuna de $i$ unidades
  ($1<i<N$) y juega hasta que pierde todo o llega a $N$.

  \smallskip
  \begin{enumerate}[label=\alph*)]
    \item Construya una cadena de Markov que describa la fortuna del jugador en cada
          instante. Incluya las probabilidades de transici\'on.
    \item El jugador al llegar a $N$ cambia su estrategia: apuesta \emph{doble o
            nada}, de modo que con probabilidad $p$ su riqueza es $2N$ (y se retira),
          mientras con probabilidad $1-p$ pierde todo. Modele esta nueva situaci\'on.
    \item Si $p=1/2$, ?`de qu\'e depende que el jugador finalmente gane o pierda?
          Sin hacer c\'alculos entregue valores espec\'ificos cuando sea posible e
          interprete sus resultados.
    \item Resuelva el problema general: encuentre las probabilidades de terminar
          ganando y perdiendo si se empieza con fortuna $i$ ($1<i<N$) y
          $p\ne 1-p$.
  \end{enumerate}
\end{enunciado}

\vspace{6pt}

% ───────────────────────────────────────────────────────────────────────
\parteheader{Parte a) \textperiodcentered\ Cadena de Markov: modelo b\'asico}

\textbf{Espacio de estados:} $\mathcal{S}=\{0,1,2,\ldots,N\}$, donde $X_n$ es la
fortuna del jugador tras $n$ jugadas. Las probabilidades de transici\'on son:
\[
  p_{k,k+1}=p, \quad p_{k,k-1}=1-p \quad (k=1,\ldots,N-1); \qquad
  p_{00}=1, \quad p_{NN}=1.
\]

\begin{center}
  \begin{tikzpicture}[
      every node/.style={circle, draw=sectionbg, fill=sectionbg!10,
          minimum size=1.0cm, font=\small\sffamily},
      every edge/.style={draw, -Stealth, thick},
      lbl/.style={draw=none, fill=none, font=\scriptsize, inner sep=2pt},
    ]
    \node (s0) at (0,   0) {$0$};
    \node (s1) at (2.2, 0) {$1$};
    \node[draw=none,fill=none,minimum size=0pt] (sd) at (4.2,0) {\large$\cdots$};
    \node (si) at (6.2, 0) {$i$};
    \node[draw=none,fill=none,minimum size=0pt] (sd2) at (8.2,0) {\large$\cdots$};
    \node (sN) at (10.2, 0) {$N$};

    \draw (s0) edge[loop above, looseness=5] node[lbl,above] {$1$} (s0);
    \draw (sN) edge[loop above, looseness=5] node[lbl,above] {$1$} (sN);

    \draw (s1) edge[bend left=20]  node[lbl,above] {$p$}     (si);
    \draw (si) edge[bend left=20]  node[lbl,below] {$1-p$}   (s1);

    \draw (s1) edge                node[lbl,below] {$1-p$}   (s0);
    \draw (si) edge[bend left=15]  node[lbl,above] {$p$}     (sN);
    \draw (sN) edge[bend left=15]  node[lbl,below] {$1-p$}   (si);
  \end{tikzpicture}
\end{center}

\textbf{Clasificaci\'on.}
\begin{itemize}
  \item $\{0\}$ y $\{N\}$: clases recurrentes absorbentes, $d=1$.
  \item $\{1,\ldots,N-1\}$: clase transitoria. Todos los estados se comunican
        (pueden ir hacia arriba y hacia abajo). No hay auto-lazos ($p_{kk}=0$),
        y todo retorno a $k$ requiere un n\'umero par de pasos ($\pm1$ deben
        compensarse). Como $\gcd(2,4,6,\ldots)=2$, el per\'iodo es $d=2$.
\end{itemize}

% ───────────────────────────────────────────────────────────────────────
\parteheader{Parte b) \textperiodcentered\ Modelo modificado: apuesta ``doble o nada'' en $N$}

Al llegar a $N$, el jugador apuesta $N$ con probabilidad $p$ de doblar ($\to 2N$)
y $1-p$ de perder todo ($\to 0$). Se agrega el estado absorbente $2N$:

\[
  p_{N,2N}=p,\quad p_{N,0}=1-p; \qquad p_{2N,2N}=1.
\]
Los dem\'as estados no cambian. El nuevo espacio de estados es
$\mathcal{S}=\{0,1,\ldots,N,2N\}$.

\textbf{Clasificaci\'on actualizada.}
\begin{itemize}
  \item $\{0\}$: recurrente absorbente (``ruina''), $d=1$.
  \item $\{2N\}$: recurrente absorbente (``gran ganancia''), $d=1$.
  \item $\{N\}$: transitorio singleton --- desde $N$ se va a $0$ o $2N$, ninguno
        de los cuales puede regresar a $N$.
  \item $\{1,\ldots,N-1\}$: clase transitoria, $d=2$ (igual que en parte a).
\end{itemize}

% ───────────────────────────────────────────────────────────────────────
\parteheader{Parte c) \textperiodcentered\ Caso $p=\tfrac{1}{2}$: ¿de qu\'e depende el resultado?}

Con $p=\tfrac{1}{2}$, la cadena es una \emph{caminata aleatoria sim\'etrica}: en
cada paso se gana o pierde una unidad con igual probabilidad. Bajo esta condici\'on,
la fortuna $X_n$ es una \textbf{martingala}:
\[
  \E{X_{n+1}\mid X_n=k} = p(k+1)+(1-p)(k-1)\big|_{p=1/2} = k.
\]
Sea $T=\min\{n:X_n\in\{0,N\}\}$ el tiempo de parada. Por el Teorema de Parada
Opcional:
\[
  \E{X_T} = X_0 = i.
\]
Como $X_T\in\{0,N\}$:
\[
  N\cdot\Prob{\text{ganar}}+0\cdot\Prob{\text{perder}}=i
  \implies
  \boxed{\Prob{\text{ganar}\mid X_0=i}=\frac{i}{N},\quad
    \Prob{\text{perder}\mid X_0=i}=\frac{N-i}{N}.}
\]

\textbf{Interpretaci\'on.} El resultado depende \emph{\'unicamente} de la posici\'on
inicial $i$ relativa al objetivo $N$: es la proporci\'on del camino ya recorrido.
Valores espec\'ificos: si $i=N/2$ (punto medio), $\Prob{\text{ganar}}=1/2$; si
$i=1$, $\Prob{\text{ganar}}=1/N$ (muy dif\'icil ganar); si $i=N-1$,
$\Prob{\text{ganar}}=(N-1)/N$ (casi seguro ganar).

% ───────────────────────────────────────────────────────────────────────
\parteheader{Parte d) \textperiodcentered\ Caso general $p\ne 1-p$: probabilidades de absorci\'on}

Sea $f_i=\Prob{\text{llegar a }N\text{ antes que a }0\mid X_0=i}$. Las condiciones
de frontera son $f_0=0$ y $f_N=1$. Para $1\le i\le N-1$, condicionando en el
primer paso:
\[
  f_i = p\,f_{i+1}+(1-p)\,f_{i-1}.
\]
Reordenando: $p\,f_{i+1}-f_i+(1-p)\,f_{i-1}=0$. Se ensaya $f_i=r^i$:
\[
  p\,r^2 - r + (1-p) = 0 \implies
  r = \frac{1\pm|1-2p|}{2p}.
\]
Para $p\ne\tfrac{1}{2}$ las dos ra\'ices son distintas:
\[
  r_1 = 1, \qquad r_2 = \frac{1-p}{p} \;=:\; \rho.
\]
La soluci\'on general es $f_i = A + B\rho^i$. Aplicando las condiciones de frontera:
\[
  f_0=0:\; A+B=0 \implies A=-B;\qquad
  f_N=1:\; -B+B\rho^N=1 \implies B=\frac{1}{\rho^N-1}.
\]
Por tanto:
\[
  f_i = \frac{\rho^i-1}{\rho^N-1} = \frac{1-\rho^i}{1-\rho^N},
  \qquad \rho=\frac{1-p}{p}\ne 1.
\]
Las probabilidades finales son:
\[
  \boxed{\Prob{\text{ganar}\mid X_0=i} = \frac{1-\rho^i}{1-\rho^N},
    \qquad
    \Prob{\text{perder}\mid X_0=i} = \frac{\rho^i-\rho^N}{1-\rho^N}.}
\]

\begin{nota}
  \textbf{Verificaci\'on:} la suma es $(1-\rho^i+\rho^i-\rho^N)/(1-\rho^N)=1$. Para
  $p>\tfrac{1}{2}$ ($\rho<1$): $\rho^N\to 0$ con $N\to\infty$, luego
  $\Prob{\text{ganar}}\to 1-\rho^i>0$ (juego favorable, probabilidad positiva de
  ganar incluso contra un adversario con recursos ilimitados). Para $p<\tfrac{1}{2}$
  ($\rho>1$): $\Prob{\text{ganar}}\to 0$ con $N\to\infty$ (juego desfavorable,
  ruina casi segura). La f\'ormula del caso $p=\tfrac{1}{2}$ ($f_i=i/N$) se
  recupera como l\'imite $\rho\to 1$ por L'H\^opital.
\end{nota}

\end{document}
